\section{Related Work}

\subsection{Federated Learning in Intelligent Transportation Systems}

The integration of Federated Learning (FL) into Intelligent Transportation Systems has gained significant momentum as a solution to the inherent privacy risks associated with centralized traffic data collection. Recent literature extensively explores the application of FL across diverse vehicular contexts, including connected and automated vehicles (CAVs), traffic flow prediction, and edge-assisted route optimization \cite{Sanadi_2026}. For instance, frameworks such as FedGRU have demonstrated the ability to accurately forecast macroscopic traffic volumes without exposing the raw spatial-temporal trajectories of individual drivers \cite{Liu_2020}. Similarly, advanced architectures leveraging decentralized spatial-temporal networks (e.g., FedSTN) have been proposed to manage urban traffic heterogeneity across disparate geographical zones \cite{Yuan_2024, Zhang_2025}. These existing federated learning frameworks have primarily focused on privacy-preserving distributed training and communication efficiency. However, these studies mainly optimize communication, privacy, or training accuracy, while the problem of maintaining reliable inference under incomplete feature availability remains insufficiently explored.

\subsection{Traffic Anomaly Detection and Trajectory Analysis}

The detection of anomalies in vehicular networks has historically relied heavily on centralized data mining and deep learning architectures. Traditional methods often process raw, high-frequency GPS trajectories using grid-based clustering techniques to identify spatial deviations in real-time \cite{Raiyn_2014}. While effective in controlled settings, these trajectory-mining approaches generally struggle with the massive distance computation complexities inherent in large-scale urban grids \cite{Su_2019}. 

To capture complex spatial-temporal correlations, contemporary research has pivoted towards deep learning models. For example, centralized ResNet models equipped with attention mechanisms have successfully addressed the challenge of imbalanced incident data, achieving high accuracy in predicting anomalies across highway datasets \cite{El_Manaa_2022}. Furthermore, Spatial-Temporal Graph Neural Networks (ST-GNNs) have been proposed to map non-Euclidean road segments, forecasting traffic conditions and classifying congestion events with remarkable precision \cite{Gao_2023}. 

Despite their predictive power, many of these trajectory-based and graph-based approaches rely on processing high-dimensional GPS data or maintaining centralized spatial representations. Operating such architectures requires significant computational overhead, rendering them highly prohibitive for direct deployment on resource-constrained edge devices (e.g., roadside units or on-board computers). Moreover, transmitting the requisite spatial data to a central server fundamentally violates V2X privacy constraints. However, these studies mainly optimize spatial prediction accuracy or graph computation efficiency, while the problem of maintaining reliable anomaly classification locally on edge devices without exposing raw trajectories remains insufficiently explored.

\subsection{Edge Deployment Challenges and Out-of-Distribution Mitigation}

The deployment of deep learning models on intelligent edge devices is frequently bottlenecked by volatile operational environments. Research into adaptive federated optimization, such as FedYogi and FedAdam, has made significant strides in mitigating statistical heterogeneity by dynamically adjusting server-side learning rates \cite{Yang_2024}. Additionally, control algorithms for local updates have been developed to balance computational resource constraints at the network edge \cite{Wang_2020}. While these methods effectively address non-IID (Independent and Identically Distributed) training data divergence \cite{Lu_2024}, they generally operate under the assumption that inference-time inputs will adhere to the trained feature distributions. 

When deploying autonomous edge sensors, physical hardware faults and connectivity drops frequently result in extreme Out-of-Distribution (OOD) feature vectors. Many OOD detection techniques in computer vision and natural language processing rely on confidence calibration, such as temperature scaling (e.g., ODIN), or post-hoc scoring methods \cite{Shalev_2022}. However, their direct applicability to resource-constrained intelligent transportation deployments remains an active area of research. These traditional methods typically intervene only after the model has processed the malformed input, which can still lead to mathematical saturation and delayed response times on edge hardware. However, these studies mainly optimize server-side model aggregation or post-hoc prediction confidence, while the problem of maintaining reliable inference through dynamic anomaly warnings and imputation under incomplete feature availability remains insufficiently explored.

\subsection{Summary of Research Gap}

Collectively, the existing literature has made substantial progress in privacy-preserving federated learning, traffic anomaly detection, and edge intelligence. However, these research directions have largely evolved independently, leaving the challenge of maintaining reliable inference under incomplete feature availability in decentralized intelligent transportation systems comparatively underexplored. This observation motivates the deployment-aware methodology presented in the next section.
